%%=============================================================================
%% CAPITULO 4 -- APRESENTACAO, ANALISE E DISCUSSAO DOS DADOS
%% Esqueleto provisorio (dissertacao dummy): texto a ser desenvolvido.
%%=============================================================================

\section{Análise e discussão dos dados}

A execução do artefato classificou as 760 videoaulas ativas do acervo,
distribuídas entre quinze disciplinas e cinco áreas de referência, sem registro
de falha de processamento. A análise do comportamento da classificação revelou
distribuição de escores unimodal e concentrada e margens reduzidas entre as
primeiras posições do ordenamento: com margens dessa ordem, a habilidade que
ocupa a primeira posição não representa escolha nitidamente destacada, mas o
primeiro elemento de um conjunto de descritores quase equidistantes do conteúdo
da aula.

A Tabela~\ref{tab:concordancia} apresenta a concordância das alternativas com a
atribuição oficial de habilidades do INEP sobre os 178 itens válidos da edição
de 2023 do ENEM, medida na primeira posição do ordenamento e nas cinco
primeiras posições, com trinta habilidades candidatas por área.

\begin{table}[htb]
  \caption{Concordância com a atribuição oficial de habilidades do INEP (ENEM 2023, N = 178)}
  \label{tab:concordancia}
  \centering
  \begin{tabular}{lrr}
    \toprule
    \textbf{Condição} & \textbf{Primeira posição (\%)} & \textbf{Cinco posições (\%)} \\
    \midrule
    Acaso (30 candidatas)                         &  3,3 & 16,7 \\
    Lexical (TF-IDF), sem inteligência artificial & 11,8 & 32,6 \\
    Embeddings BGE-M3 (método do artefato)        & 15,7 & 43,8 \\
    Modelo generativo proprietário                & 41,0 & 78,1 \\
    \bottomrule
  \end{tabular}
  \fonte{Elaborado pelo autor (2026), a partir dos resultados da pesquisa.}
\end{table}

Três leituras decorrem da tabela. Todas as condições superam o acaso com folga
suficiente para que a medição não seja ruído. O artefato supera o método sem
inteligência artificial, 15,7\% contra 11,8\% na primeira posição e 43,8\%
contra 32,6\% em cinco posições, o que atende ao parâmetro de aceitação fixado
a priori, com a ressalva de que, em Linguagens e Códigos, o TF-IDF supera os
embeddings na primeira posição. E o artefato alcança menos da metade do
desempenho da condição generativa, 15,7\% contra 41,0\%, o que mostra que o
patamar observado não decorre de limite intrínseco da tarefa.

Para uma rede pública, portanto, a decisão não é qual método é melhor, e sim
qual combinação de qualidade, custo e autonomia é aceitável. A condição
generativa classifica melhor, mas transmite conteúdo educacional público a
infraestrutura de terceiro e depende de versão mantida por fornecedor externo;
as condições locais executam em estação de trabalho de prateleira, sem custo
por requisição e com reprodutibilidade integral. A comparação estabelece ainda
que o texto com que a matriz enuncia a habilidade é representação menos fiel do
que o exame cobra do que o conjunto dos itens oficialmente atribuídos a essa
habilidade.
